\lecture{03}{Структура курса}

\sourcevideo
    {https://www.youtube.com/playlist?list=PLOzRYVm0a65fhKDS8cYIC7YCuVGJ4FOJx}
    {Week 1: Lecture 3: Course Outline}

Курс объединяет выпуклую оптимизацию, анализ динамических систем,
ускоренные методы и распределённые алгоритмы. Основная часть
результатов формулируется для выпуклых задач. Отдельно будут
рассматриваться некоторые невыпуклые функции, удовлетворяющие
неравенству Поляка--Лоясиевича.

\section{Почему основное внимание уделяется выпуклым задачам}

Выпуклая функция не обязана иметь единственный минимум. Например,
постоянная функция является выпуклой, и каждая точка её области
определения является глобальным минимумом. Аналогично минимум может
достигаться на целом отрезке или более общем выпуклом множестве.

Ключевое отличие выпуклой задачи от общей невыпуклой задачи состоит
не в единственности решения, а в отсутствии неглобальных локальных
минимумов.

\begin{theorem}[Локальная и глобальная оптимальность]
Пусть функция \(f\) выпукла на выпуклом множестве \(X\). Тогда всякая
точка локального минимума функции \(f\) на \(X\) является точкой её
глобального минимума на \(X\).
\end{theorem}

Таким образом, если алгоритм нашёл локальный минимум выпуклой задачи,
дополнительная проверка его глобальности не требуется. Различные
точки минимума могут существовать, но значение целевой функции во
всех них одинаково.

В невыпуклой задаче функция может иметь несколько локальных минимумов
с различными значениями:
\[
    f(x_{\mathrm{loc}})
    >
    f(x^\star).
\]
Поэтому локальная сходимость алгоритма в общем случае не гарантирует
нахождение глобального решения.

\begin{remark}
Формальные определения выпуклых множеств, выпуклых функций и строгой
выпуклости вводятся в последующих лекциях. Здесь используется только
главное следствие выпуклости: локальный минимум является глобальным.
\end{remark}

\section{Выпуклая оптимизация}

Первая содержательная часть курса посвящена централизованным задачам
оптимизации. Будут рассмотрены как задачи без ограничений,
\[
    \min_{x\in\RR^d} f(x),
\]
так и задачи с ограничениями,
\[
    \begin{aligned}
        \min_{x\in\RR^d}
        \quad& f(x),
        \\
        \text{при условии}
        \quad& x\in\mathcal X.
    \end{aligned}
\]

Основной класс составляют выпуклые функции и выпуклые допустимые
множества. Помимо них рассматриваются функции, удовлетворяющие
неравенству Поляка--Лоясиевича. Такие функции могут быть невыпуклыми,
но сохранять свойства, достаточные для анализа быстрой сходимости.

В качестве мотивационного примера используется функция
\[
    f(x)=x^2+3\sin^2x.
\]
Она не является выпуклой на всей прямой, но имеет единственный
глобальный минимум в точке \(x^\star=0\). Этот пример показывает, что
невыпуклость не исключает простой глобальной структуры задачи, однако
сама по себе не гарантирует выполнения PL-неравенства.

\section{Теория устойчивости}

После основ выпуклой оптимизации курс переходит к теории устойчивости
по Ляпунову для непрерывных динамических систем
\[
    \dot{x}=F(x).
\]

Оптимизационный метод сначала представляется как динамическая
система, множество равновесий которой должно совпадать с множеством
решений исходной задачи. Затем устойчивость равновесий используется
для доказательства сходимости, а её тип задаёт скорость приближения к
решению.

Особое внимание уделяется фиксированно-временной сходимости.

\begin{definition}[Фиксированно-временная сходимость]
Равновесие \(x^\star\) обладает фиксированно-временной сходимостью,
если существует число \(T_{\max}>0\), не зависящее от начального
состояния, такое что всякая траектория достигает равновесия не позднее
момента \(T_{\max}\):
\[
    x(t)=x^\star
    \qquad
    \text{при всех }t\ge T(x_0),
    \qquad
    T(x_0)\le T_{\max}.
\]
\end{definition}

Главная особенность этой оценки состоит в независимости верхней
границы времени сходимости от начальной точки. Это сильнее обычной
асимптотической или экспоненциальной сходимости, при которой решение
приближается к равновесию, но в общем случае не достигает его за
конечное время.

\section{Проектирование ускоренных методов}

Непрерывные динамические системы используются не только для анализа
известных алгоритмов, но и для построения новых методов. Общая схема
имеет вид
\[
    \begin{gathered}
        \text{оптимизационная задача}
        \longrightarrow
        \text{непрерывная динамика},
        \\
        \text{непрерывная динамика}
        \longrightarrow
        \text{анализ устойчивости},
        \\
        \text{анализ устойчивости}
        \longrightarrow
        \text{дискретный алгоритм}.
    \end{gathered}
\]

Сначала выбирается система
\[
    \dot{x}=F(x),
\]
множество равновесий которой совпадает с множеством решений, а скорость
сходимости удовлетворяет требуемой оценке. После этого система
дискретизируется:
\[
    x^{k+1}
    =
    x^k+\eta_kF(x^k).
\]

Непрерывная система может иметь фиксированно-временную или иную
усиленную сходимость, однако после дискретизации соответствующие
гарантии необходимо исследовать заново. Поведение дискретного метода
зависит от выбранной схемы и величины шага \(\eta_k\).

\section{Переход к распределённой оптимизации}

После централизованной оптимизации и теории устойчивости начинается
распределённая часть курса. Для её построения вводятся основные
понятия теории графов:
\[
    \Graph=(\Vertices,\Edges).
\]

Вершины графа соответствуют вычислительным агентам, а рёбра задают
допустимые каналы коммуникации. Будут рассмотрены связность графа и
матрица Лапласа \(L=D-A\), где \(D\) --- матрица степеней, а \(A\) ---
матрица смежности.

\begin{definition}[Значение Фидлера]
Для неориентированного графа значением Фидлера называется второе
наименьшее собственное значение его матрицы Лапласа:
\[
    \lambda_2(L).
\]
\end{definition}

Значение Фидлера называют также алгебраической связностью. Для
неориентированного графа условие \(\lambda_2(L)>0\) эквивалентно его
связности.

После введения графовой модели рассматриваются алгоритмы, в которых
каждый агент выполняет локальные вычисления и обменивается информацией
только с соседями. Их цель состоит в решении глобальной задачи без
полного доступа к данным и функциям остальных участников.

Таким образом, содержательная последовательность курса имеет вид
\[
    \begin{gathered}
        \text{выпуклая оптимизация}
        \longrightarrow
        \text{теория устойчивости},
        \\
        \text{теория устойчивости}
        \longrightarrow
        \text{ускоренные методы},
        \\
        \text{теория графов}
        \longrightarrow
        \text{консенсус},
        \\
        \text{консенсус}
        \longrightarrow
        \text{распределённая оптимизация}.
    \end{gathered}
\]

Заключительная часть курса посвящена крупномасштабному и
федеративному обучению.

\section{Источники курса}

Курс не следует одному учебнику. Материал основывается на конспектах
лекций и исследовательских статьях, сопровождающих соответствующие
темы.

Значительная часть рассматриваемых результатов относится к
современным направлениям ускоренной, крупномасштабной и
распределённой оптимизации. Эти методы образуют теоретическую основу
распределённого машинного обучения.


Основу курса составляют выпуклая оптимизация и функции, допускающие
усиленные оценки сходимости. Затем вводится теория устойчивости
непрерывных динамических систем, включая сходимость за фиксированное время, и на её основе проектируются ускоренные методы. Далее
изучаются теория графов, консенсус, распределённая оптимизация,
крупномасштабное и федеративное обучение.